\documentclass[letterpaper]{article}
\usepackage{aaai}

\usepackage{times}
\usepackage{helvet}
\usepackage{courier}

\usepackage{graphicx}

\pdfinfo{
/Title (YARP)
/Subject (YARP)
/Author (IIT)}

\title{YARP, a Thin Middleware for (Humanoid) Robots}
\author{Paul Fitzpatrick{\rm ,}  Lorenzo Natale \and Giorgio Metta \\
Italian Institute of Technology \\
Via Morego, 30 \\16163 Genova, Italy}
\setcounter{secnumdepth}{0}

\begin{document} 
\maketitle
\begin{abstract}
\begin{quote}

YARP stands for ``Yet Another Robot Platform.''  It is a robot
middleware that began life as a thin layer over the QNX real-time
operating system to adapt it for use by humanoid robots.  It is now
used on all kinds of operating systems and robots around the world.
It is free and open software, released under the LGPL.  Over the past
decade, YARP's communication model has proven to be at a sweet spot
that combines efficiency, flexibility, and ease of use.

\end{quote}
\end{abstract}

\noindent 
YARP is used to build robot control systems as a collection of
programs communicating in a peer-to-peer way, using an open-ended
family of connection types (TCP, UDP, multicast, local, shared memory, 
MPI, mjpeg-over-http, XML/RPC, TCPROS, plain-text, etc.) that
can be swapped in and out
as need dictates.
It also supports similarly flexible interfacing
with hardware devices.  The strategic goal
of YARP's developers is to increase the longevity
of robot software projects \cite{fitzpatrick08towards}.

YARP's first version coalesced in 2000.  It was 
shaped by the problem of doing productive research despite constant flux
in our robot platforms (hence the name).  Since then, experience with
incompatible architectures, frameworks, and middleware~-- which we like
to call collectively ``muddleware''~-- has taught us to make YARP a {\it
  reluctant} middleware, with no desire or expectation to be in
control of a user's system.  Some see YARP's restraint as a
negative, but for others it is appealing; a typical user comment in a
recent survey was:  ``It's lightweight and easy to use. I like the
multi-platform support, and its non-monopolistic philosophy.''

\begin{figure}
\centerline{\includegraphics[width=5cm]{icub.jpg}}
\caption{YARP works on
complex humanoids such as the iCub, on embedded systems, and 
on everything in between.
} 
\end{figure}

Communication in YARP generally follows the Observer design pattern
\cite{gamma95design}. {\tt Port} objects deliver messages to any
number of observers (other {\tt Port}s), in any number of processes,
distributed across any number of machines, using any of many
underlying communication protocols.  {\tt Port}s can be connected on
an individual basis or as topic-based groups.  A message from a single
{\tt Port} may be sent simultaneously across multiple connections
using distinct protocols.  {\tt Port}s can be connected to non-YARP
network entities, such as IP cameras, web servers, or ROS nodes
\cite{quigley2009ros}.
All of this can require some fancy footwork by YARP, such as
grouping multicast connections, or coordinating with a foreign
middleware.  YARP has a similarly flexible device interface.  
Individual YARP-using programs can be upgraded
over time to deal with changes in hardware or networking without
touching user source code, and without needing to make a ``big bang''
change of all programs at once.


\section{Stay out of my build}

YARP was born in 2000 on an early humanoid robot (called Kismet)
controlled by a set of Motorola 68332 processors, an Apple Mac, and a
loose network of PCs running QNX, Linux, and Microsoft Windows
\cite{metta06yarp}.  Communication on this robot was an ad-hoc
mixture of dual-port RAM polling, QNX message passing, CORBA, and raw
sockets.  From the beginning, YARP has been built with the implicit
assumption that it is just one part of the users' environment.  So
YARP restrains itself to behave like any other library, such as
OpenCV, ODE, etc., rather than expecting any special treatment.

But there's a good reason why most middleware are more than just
regular libraries: marshalling/demarshalling.  To send
user data between programs, that data must be translated to and from
its format in memory to its format ``on the
wire.''  It is commonplace to write a parser to convert a user's
structure description into generated code for performing this
conversion.
YARP has historically skirted around this approach, for two main
reasons.  First, it is invasive, since it insinuates the middleware
into the user's build toolchain (imagine how that much fun trying to
use {\it two} such middleware is).  Second, it can be inefficient,
potentially leading to time wasted making copies of large data structures
(YARP goes to great lengths to make sure that large structures such as
images get transmitted with zero extra copies made).  As computer
scientists, it is hard to resist adding automation for this, 
but the cost of meeting that automation's requirements should be borne 
in mind.  A quote from a user:

\begin{quote}
``Compared to, e.g., CAVIAR and Psyclone, YARP
looks like a fairly standard library~- neither does it
do its own message scheduling nor does it provide
heavy-handed semantics for message definitions or
networking. That may be its very strength.'' \cite{stefansson09yarp}
\end{quote}


\section{The {\tt telnet} test}

We've found that a useful way to evaluate a middleware is
this: can a user monitor and insert traffic between programs using just a
{\tt telnet} client?  The {\tt telnet} program can open (more or less)
raw TCP sessions, and is commonly used to manually send messages to
servers speaking SMTP, HTTP, IRC etc.  
We started applying the {\tt telnet} test to YARP out of past frustration with other
middleware, where simply passing a few numbers to a collaborator's
program required jumping through a dozen hoops.  We therefore
(following the ``golden rule'') make it easy for others
to send a few numbers to a YARP-based program without having to dig
through protocol specifications, or link against our libraries, or use
our build machinery~-- all of which can be a time-sink.
This is similar to Google's ``data liberation front,''
a subproject to ensure that it is easy for users to migrate
their data from a service, avoiding
lock-in.

A YARP network is designed to be usable without YARP.  This doesn't
happen by accident.  
For example, we took care that YARP
connections could be initiated by either the ``sender'' or ``receiver,''
with the logical flow of data being freely reversable.  This is
important for supporting a wide range of protocols, which may be {\it
  pull} or {\it push} in nature, but it is also means a foreign
program can both send input to and read from a YARP program without
getting stuck writing a server for at least
one of the directions of data flow.
Examples of making YARP connections without using the YARP libraries are
available in C, Python, and Tcl, and of course for {\tt telnet}.
Users seem to value the interoperability that YARP provides:

\begin{quote}
``YARP was chosen as the communication library with
which all communication protocols were implemented as
one of the goals of the design of the communication stack
was to make it possible to interact with programs that are
developed without using MeRMaID.'' \cite{barbosa09mermaid}
\end{quote}

\section{Conclusions}

Why {\it not} use YARP?  
YARP's commitment to portability slows its growth, since taking on any
new dependency is complicated.  
A full native implementation of YARP
exists only in C++; we rely on SWIG for wrappers in other languages.
Lack of an interface definition
language (IDL) and associated code generation can lead to some tedium
implementing classic RPC-style code. YARP currently uses
a central name server to do match-making, which can be 
problematic for applications such as modular
robotics where there is no clear central hub\footnote{{\tt LCM}
\cite{huang2010lcm} has an interesting
solution to this: it commits to UDP multicast for all messages,
with broadcast messages filtered by clients; in this case, 
a central name server is not needed, reducing a point of
failure.}.
The LGPL license that YARP is
under is commercial-use friendly, but still could
complicate certain proprietary uses of the library.

Why use YARP?  If you're looking for something portable, light-weight,
and flexible, then YARP is it.  YARP serves a truly diverse,
interdisciplinary community.  While YARP grew up on
high-end humanoids with lots of resources, it works fine on embedded
systems.  It has a flexible and open model of connections that has
stood the test of time.  It places no constraints on the user's build
system (though we are big fans of CMake).  And it passes the ``{\tt
  telnet} test,'' a property of just about
every successful, durable network protocol or
format.

\bibliography{yarp}
\bibliographystyle{aaai}

\end{document}
